% ======================================================================
% SECTION 14: AUTONOMOUS MEDICAL WRITING AND DISCLOSURE
% File: sections/writing_disclosure.tex
% ======================================================================

Medical writing and disclosure complete the autonomous sponsor's document lifecycle, from
protocol generation at trial initiation through archival 25 years after trial completion.
The writing and disclosure agent automates the generation of all sponsor-authored documents
required by regulators, ethics committees, trial registries, and the scientific community.
This section covers the document generation pipeline, regulatory disclosure requirements,
and long-term archiving infrastructure.

\subsection{Writing and Disclosure Agent}

The writing and disclosure agent (writing\_disclosure\_agent) replaces the medical writing
and publications team. It generates clinical trial documents from structured data produced
by other agents, ensuring consistency across all documents and compliance with applicable
formatting standards.

Protocol document generation produces ICH M11-compliant structured protocols \cite{ich-m11}
from the specifications created by the clinical accountability agent and the biostatistics
agent. The writing agent formats the protocol with appropriate section structure, cross-references,
and appendices. Protocol amendments are generated with integrated change tracking that
highlights modifications relative to the previous version and documents the rationale for
each change.

Investigator's Brochure authoring produces the comprehensive reference document required
under 21 CFR \S312.55 \cite{cfr-part-312}. The agent compiles nonclinical pharmacology and
toxicology data, clinical pharmacology summaries, prior clinical experience, and the current
safety profile into the structured IB format. IB updates are generated when new safety
information, clinical data, or manufacturing changes warrant revision, with version tracking
and distribution documentation.

Clinical Study Report generation follows the ICH E3 structure \cite{ich-e3} for the
comprehensive report of a clinical trial. The agent produces all CSR sections from
study-level data: synopsis, trial conduct summary, efficacy and safety results, statistical
analysis outputs (tables, listings, and figures), individual patient data listings, and
appendices. The CSR integrates statistical outputs from the biostatistics agent, safety
narratives from the safety agent, and regulatory compliance documentation from the quality
agent.

Table~\ref{tab:document-pipeline} presents the document generation pipeline.

\begin{longtable}{L{2.0cm} L{2.2cm} L{2.5cm} L{2.0cm} L{2.8cm}}
\caption{Document Generation Pipeline}
\label{tab:document-pipeline} \\
\toprule
\textbf{Document Type} & \textbf{Regulatory Basis} & \textbf{Generation Trigger} & \textbf{Agent Owner} & \textbf{Review Gate} \\
\midrule
\endfirsthead
\toprule
\textbf{Document Type} & \textbf{Regulatory Basis} & \textbf{Generation Trigger} & \textbf{Agent Owner} & \textbf{Review Gate} \\
\midrule
\endhead
\midrule
\multicolumn{5}{r}{\textit{Continued on next page}} \\
\endfoot
\bottomrule
\endlastfoot
Protocol & ICH M11, ICH E6(R3) & Protocol design completion by clinical accountability agent & writing\_disclosure\_agent & Sponsor-of-record approval \\
Protocol amendment & ICH E6(R3), 21 CFR 312.30 & Protocol change authorized by clinical accountability agent & writing\_disclosure\_agent & Sponsor-of-record approval \\
Investigator's Brochure & 21 CFR 312.55, ICH E6(R3) & Initial IND filing or new safety data & writing\_disclosure\_agent & Medical monitor review \\
Informed consent form & 21 CFR Part 50, adapted for Physical AI & Protocol finalization, site-specific customization & writing\_disclosure\_agent & IRB approval required \\
IND annual report & 21 CFR 312.33 & Annual anniversary of IND effective date & regulatory\_agent & Sponsor-of-record review \\
DSUR & ICH E2F & Annual (within 60 days of DIBD) & safety\_agent & Medical monitor approval \\
CSR & ICH E3 & Database lock and analysis completion & writing\_disclosure\_agent & Sponsor-of-record approval \\
Statistical tables/listings & ICH E3, SAP & Analysis dataset availability & data\_biostats\_agent & Biostatistics QC \\
\end{longtable}

\subsection{Regulatory Document Generation}

Regulatory document generation extends beyond the core clinical documents to encompass
the complete set of sponsor-authored regulatory documents. IND module assembly produces
the five eCTD modules from component documents generated by specialized agents: Module 1
(administrative information, cover letters, application forms), Module 2 (summaries of
quality, nonclinical, and clinical data), Module 3 (quality/CMC data from the supply
agent), Module 4 (nonclinical study reports), and Module 5 (clinical study reports and
individual patient data). The regulatory agent assembles these modules into eCTD v4.0
\cite{fda-ectd} compliant packages.

Annual report generation compiles 21 CFR \S312.33 required content: brief summary of
the study status, summary of adverse experiences, IND safety reports filed during the
year, literature references, manufacturing changes, and the current Investigator's
Brochure. The agent generates the annual report automatically from cumulative study data,
safety database exports, and regulatory submission logs.

Informed consent document generation produces the informed consent form and any supplementary
information sheets required for patient enrollment. For Physical AI trials, the consent
documents include the additional elements required by the adapted 21 CFR Part 50
\cite{cfr50-adapt}: description of the robotic systems involved, explanation of the safety
monitoring specific to robotic procedures, description of the human oversight
mechanisms, and the participant's right to request human-only alternatives where available.

\subsection{Disclosure and Publication Management}

Disclosure management automates compliance with the multiple, overlapping disclosure
requirements that apply to clinical trials. Table~\ref{tab:disclosure-timeline} presents
the disclosure timeline requirements across major registries.

\begin{longtable}{L{2.0cm} L{2.2cm} L{2.0cm} L{3.0cm} L{2.4cm}}
\caption{Disclosure Timeline Requirements}
\label{tab:disclosure-timeline} \\
\toprule
\textbf{Registry} & \textbf{Regulation} & \textbf{Deadline} & \textbf{Content} & \textbf{Agent Automation} \\
\midrule
\endfirsthead
\toprule
\textbf{Registry} & \textbf{Regulation} & \textbf{Deadline} & \textbf{Content} & \textbf{Agent Automation} \\
\midrule
\endhead
\midrule
\multicolumn{5}{r}{\textit{Continued on next page}} \\
\endfoot
\bottomrule
\endlastfoot
ClinicalTrials.gov & FDAAA 801 & 1 year from primary completion & Registration, results summary, adverse events, protocol, SAP & Automated results extraction and posting \\
ClinicalTrials.gov & 42 CFR Part 11 & 30 days from study initiation & Protocol registration with required data elements & Automated registration from protocol \\
CTIS (EU) & EU CTR 536/2014 & 1 year from completion (6 months pediatric) & Results summary, lay summary & Automated extraction with lay summary generation \\
CTIS (EU) & EU CTR 536/2014 & 30 days post-MA decision & CSR summary & Automated CSR summary extraction \\
EudraCT & EU CTR (legacy) & Per transition timeline & Results from pre-CTIS trials & Migration support from legacy format \\
WHO ICTRP & WHO policy & Registration before first enrollment & ICTRP minimum dataset & Automated registration from protocol \\
\end{longtable}

ClinicalTrials.gov results posting follows FDAAA 801 \cite{fdaaa-801} and 42 CFR Part 11
requirements. The agent extracts results data from the statistical analysis outputs,
formats them in the ClinicalTrials.gov results data element structure, generates adverse
event tables in the required format, and prepares the results for submission. A quality
review step verifies completeness and consistency before posting. The one-year deadline
from primary completion date is tracked by the compliance calendar managed by the
regulatory agent.

CTIS results and lay summary posting follows the EU Clinical Trials Regulation
requirements. Results summaries must be posted within one year of trial completion
(six months for pediatric trials). The lay summary requires plain-language description
of trial results accessible to non-specialist readers. The writing agent generates lay
summaries using controlled vocabulary and readability assessment to ensure accessibility.

\subsection{Archiving and Long-Term Retention}

The EU Clinical Trials Regulation imposes a 25-year archiving requirement
\cite{eu-ctr} on sponsors, with designated sponsor-appointed archive persons responsible
for maintaining access to archived records throughout the retention period. The autonomous
sponsor's archiving system addresses this requirement through digital preservation
strategies that account for technology evolution over the 25-year retention window.

Essential document lifecycle management tracks every document classified as essential
under ICH E6(R3) from creation through archival. The lifecycle includes creation
(document generated by the writing agent or received from external sources), review
and approval (quality control and sponsor-of-record review), distribution (delivery to
investigators, IRBs, and regulatory authorities), maintenance (version control and
amendment tracking), and archival (transfer to long-term archive with integrity
verification).

Digital preservation implements three strategies to ensure long-term accessibility:
format standardization (all documents archived in open formats: PDF/A-3 for documents
with embedded attachments, XML for structured data, CSV for tabular datasets), media
migration (periodic transfer to current storage media before obsolescence), and
integrity verification (cryptographic hash validation at defined intervals to detect
data degradation). The archive system generates annual preservation status reports
documenting format validity, media condition, and integrity verification results.
